% Options for packages loaded elsewhere
\PassOptionsToPackage{unicode}{hyperref}
\PassOptionsToPackage{hyphens}{url}
%
\documentclass[
  ignorenonframetext,
]{beamer}
\usepackage{pgfpages}
\setbeamertemplate{caption}[numbered]
\setbeamertemplate{caption label separator}{: }
\setbeamercolor{caption name}{fg=normal text.fg}
\beamertemplatenavigationsymbolsempty
% Prevent slide breaks in the middle of a paragraph
\widowpenalties 1 10000
\raggedbottom
\setbeamertemplate{part page}{
  \centering
  \begin{beamercolorbox}[sep=16pt,center]{part title}
    \usebeamerfont{part title}\insertpart\par
  \end{beamercolorbox}
}
\setbeamertemplate{section page}{
  \centering
  \begin{beamercolorbox}[sep=12pt,center]{part title}
    \usebeamerfont{section title}\insertsection\par
  \end{beamercolorbox}
}
\setbeamertemplate{subsection page}{
  \centering
  \begin{beamercolorbox}[sep=8pt,center]{part title}
    \usebeamerfont{subsection title}\insertsubsection\par
  \end{beamercolorbox}
}
\AtBeginPart{
  \frame{\partpage}
}
\AtBeginSection{
  \ifbibliography
  \else
    \frame{\sectionpage}
  \fi
}
\AtBeginSubsection{
  \frame{\subsectionpage}
}
\usepackage{amsmath,amssymb}
\usepackage{lmodern}
\usepackage{iftex}
\ifPDFTeX
  \usepackage[T1]{fontenc}
  \usepackage[utf8]{inputenc}
  \usepackage{textcomp} % provide euro and other symbols
\else % if luatex or xetex
  \usepackage{unicode-math}
  \defaultfontfeatures{Scale=MatchLowercase}
  \defaultfontfeatures[\rmfamily]{Ligatures=TeX,Scale=1}
\fi
\usetheme[]{Montpellier}
% Use upquote if available, for straight quotes in verbatim environments
\IfFileExists{upquote.sty}{\usepackage{upquote}}{}
\IfFileExists{microtype.sty}{% use microtype if available
  \usepackage[]{microtype}
  \UseMicrotypeSet[protrusion]{basicmath} % disable protrusion for tt fonts
}{}
\makeatletter
\@ifundefined{KOMAClassName}{% if non-KOMA class
  \IfFileExists{parskip.sty}{%
    \usepackage{parskip}
  }{% else
    \setlength{\parindent}{0pt}
    \setlength{\parskip}{6pt plus 2pt minus 1pt}}
}{% if KOMA class
  \KOMAoptions{parskip=half}}
\makeatother
\usepackage{xcolor}
\IfFileExists{xurl.sty}{\usepackage{xurl}}{} % add URL line breaks if available
\IfFileExists{bookmark.sty}{\usepackage{bookmark}}{\usepackage{hyperref}}
\hypersetup{
  hidelinks,
  pdfcreator={LaTeX via pandoc}}
\urlstyle{same} % disable monospaced font for URLs
\newif\ifbibliography
\usepackage{color}
\usepackage{fancyvrb}
\newcommand{\VerbBar}{|}
\newcommand{\VERB}{\Verb[commandchars=\\\{\}]}
\DefineVerbatimEnvironment{Highlighting}{Verbatim}{commandchars=\\\{\}}
% Add ',fontsize=\small' for more characters per line
\usepackage{framed}
\definecolor{shadecolor}{RGB}{248,248,248}
\newenvironment{Shaded}{\begin{snugshade}}{\end{snugshade}}
\newcommand{\AlertTok}[1]{\textcolor[rgb]{0.94,0.16,0.16}{#1}}
\newcommand{\AnnotationTok}[1]{\textcolor[rgb]{0.56,0.35,0.01}{\textbf{\textit{#1}}}}
\newcommand{\AttributeTok}[1]{\textcolor[rgb]{0.77,0.63,0.00}{#1}}
\newcommand{\BaseNTok}[1]{\textcolor[rgb]{0.00,0.00,0.81}{#1}}
\newcommand{\BuiltInTok}[1]{#1}
\newcommand{\CharTok}[1]{\textcolor[rgb]{0.31,0.60,0.02}{#1}}
\newcommand{\CommentTok}[1]{\textcolor[rgb]{0.56,0.35,0.01}{\textit{#1}}}
\newcommand{\CommentVarTok}[1]{\textcolor[rgb]{0.56,0.35,0.01}{\textbf{\textit{#1}}}}
\newcommand{\ConstantTok}[1]{\textcolor[rgb]{0.00,0.00,0.00}{#1}}
\newcommand{\ControlFlowTok}[1]{\textcolor[rgb]{0.13,0.29,0.53}{\textbf{#1}}}
\newcommand{\DataTypeTok}[1]{\textcolor[rgb]{0.13,0.29,0.53}{#1}}
\newcommand{\DecValTok}[1]{\textcolor[rgb]{0.00,0.00,0.81}{#1}}
\newcommand{\DocumentationTok}[1]{\textcolor[rgb]{0.56,0.35,0.01}{\textbf{\textit{#1}}}}
\newcommand{\ErrorTok}[1]{\textcolor[rgb]{0.64,0.00,0.00}{\textbf{#1}}}
\newcommand{\ExtensionTok}[1]{#1}
\newcommand{\FloatTok}[1]{\textcolor[rgb]{0.00,0.00,0.81}{#1}}
\newcommand{\FunctionTok}[1]{\textcolor[rgb]{0.00,0.00,0.00}{#1}}
\newcommand{\ImportTok}[1]{#1}
\newcommand{\InformationTok}[1]{\textcolor[rgb]{0.56,0.35,0.01}{\textbf{\textit{#1}}}}
\newcommand{\KeywordTok}[1]{\textcolor[rgb]{0.13,0.29,0.53}{\textbf{#1}}}
\newcommand{\NormalTok}[1]{#1}
\newcommand{\OperatorTok}[1]{\textcolor[rgb]{0.81,0.36,0.00}{\textbf{#1}}}
\newcommand{\OtherTok}[1]{\textcolor[rgb]{0.56,0.35,0.01}{#1}}
\newcommand{\PreprocessorTok}[1]{\textcolor[rgb]{0.56,0.35,0.01}{\textit{#1}}}
\newcommand{\RegionMarkerTok}[1]{#1}
\newcommand{\SpecialCharTok}[1]{\textcolor[rgb]{0.00,0.00,0.00}{#1}}
\newcommand{\SpecialStringTok}[1]{\textcolor[rgb]{0.31,0.60,0.02}{#1}}
\newcommand{\StringTok}[1]{\textcolor[rgb]{0.31,0.60,0.02}{#1}}
\newcommand{\VariableTok}[1]{\textcolor[rgb]{0.00,0.00,0.00}{#1}}
\newcommand{\VerbatimStringTok}[1]{\textcolor[rgb]{0.31,0.60,0.02}{#1}}
\newcommand{\WarningTok}[1]{\textcolor[rgb]{0.56,0.35,0.01}{\textbf{\textit{#1}}}}
\setlength{\emergencystretch}{3em} % prevent overfull lines
\providecommand{\tightlist}{%
  \setlength{\itemsep}{0pt}\setlength{\parskip}{0pt}}
\setcounter{secnumdepth}{-\maxdimen} % remove section numbering
%\documentclass[unicode,10pt]{beamer}% 'unicode'が必要
\usepackage{luatexja}% 日本語を使う場合は必須
%\usepackage[japanese]{babel}	
%\usefonttheme[onlymath]{serif}
%\usepackage[T1]{fontenc} %これを使うと、ギリシャ文字が出ない
%\usepackage{textcomp}
%\usepackage[scale=1.0]{tgheros} % Sans serif font
%\usepackage[scaled]{beramono} % Monospace font
%\usepackage{luatexja-otf}
%\renewcommand{\kanjifamilydefault}{\gtdefault}
%\usepackage[haranoaji]{luatexja-preset}

% スライド番号の表示 %%%%%%%%%%%%%%%%%%%
\makeatletter
\setbeamertemplate{footline}{%
  \leavevmode%
  \hbox{%
  \begin{beamercolorbox}[wd=.5\paperwidth,ht=2.25ex,dp=1ex,right]{date in head/foot}%
    \insertframenumber{} \hspace*{2ex} % new version without total frames
  \end{beamercolorbox}}%
  \vskip0pt%
}
\makeatother
% スライド番号の表示 %%%%%%%%%%%%%%%%%%%
\ifLuaTeX
  \usepackage{selnolig}  % disable illegal ligatures
\fi

\title{第2章: 因果関係\\
(2.1. 労働市場における人種差別)}
\subtitle{今井耕介 著\\
『社会科学のためのデータ分析入門 (QSS)』}
\author{}
\date{\vspace{-2.5em}2026-03-09}

\begin{document}
\frame{\titlepage}

\hypertarget{ux52b4ux50cdux5e02ux5834ux306bux304aux3051ux308bux4ebaux7a2eux5deeux5225}{%
\section{2.1
労働市場における人種差別}\label{ux52b4ux50cdux5e02ux5834ux306bux304aux3051ux308bux4ebaux7a2eux5deeux5225}}

\begin{frame}{架空履歴書実験}
\protect\hypertarget{ux67b6ux7a7aux5c65ux6b74ux66f8ux5b9fux9a13}{}
\begin{itemize}
\tightlist
\item
  \textbf{問い}: 労働市場において、人種に基づく差別は存在するか？
\item
  \textbf{実験デザイン}:

  \begin{itemize}
  \tightlist
  \item
    架空の履歴書を多数作成し、実際の求人に応募する。
  \item
    履歴書の名前を「白人的な名前」と「黒人的な名前」に\textbf{無作為（ランダム）に割り当てる}。
  \item
    その他の経歴（学歴や経験など）は同一。
  \end{itemize}
\item
  \textbf{結果変数}:
  雇用主から面接の連絡（コールバック）があったかどうか。
\end{itemize}
\end{frame}

\hypertarget{ux30c7ux30fcux30bfux306eux8aadux307fux8fbcux307fux3068ux78baux8a8d}{%
\section{2.1.1
データの読み込みと確認}\label{ux30c7ux30fcux30bfux306eux8aadux307fux8fbcux307fux3068ux78baux8a8d}}

\begin{frame}[fragile]{resume データの読み込み (1)}
\protect\hypertarget{resume-ux30c7ux30fcux30bfux306eux8aadux307fux8fbcux307f-1}{}
\begin{itemize}
\tightlist
\item
  実験で収集されたデータをRに読み込みます。
\end{itemize}

\scriptsize

\begin{Shaded}
\begin{Highlighting}[]
\CommentTok{\# 1. ローカルに保存したresumeデータの読み込み（推奨）}
\NormalTok{resume }\OtherTok{\textless{}{-}} \FunctionTok{read.csv}\NormalTok{(}\StringTok{"resume.csv"}\NormalTok{)}

\CommentTok{\# （参考）URLから直接読み込むことも可能}
\CommentTok{\# resume \textless{}{-} read.csv("https://ayumu{-}tanaka.github.io/QSS/QSS\_Data/resume.csv")}

\CommentTok{\# 2. データの次元（行数と列数）を確認}
\FunctionTok{dim}\NormalTok{(resume)}
\end{Highlighting}
\end{Shaded}

\begin{verbatim}
## [1] 4870    4
\end{verbatim}

\normalsize
\end{frame}

\begin{frame}[fragile]{resume データの読み込み (2)}
\protect\hypertarget{resume-ux30c7ux30fcux30bfux306eux8aadux307fux8fbcux307f-2}{}
\begin{itemize}
\tightlist
\item
  データの最初の数行を表示して、変数の並びや中身を確認します。
\end{itemize}

\scriptsize

\begin{Shaded}
\begin{Highlighting}[]
\CommentTok{\# 最初の3行を表示}
\FunctionTok{head}\NormalTok{(resume, }\AttributeTok{n =} \DecValTok{3}\NormalTok{)}
\end{Highlighting}
\end{Shaded}

\begin{verbatim}
##   firstname    sex  race call
## 1   Allison female white    0
## 2   Kristen female white    0
## 3   Lakisha female black    0
\end{verbatim}

\normalsize
\end{frame}

\hypertarget{ux30afux30edux30b9ux96c6ux8a08ux8868}{%
\section{2.1.2
クロス集計表}\label{ux30afux30edux30b9ux96c6ux8a08ux8868}}

\begin{frame}[fragile]{人種とコールバックの関係}
\protect\hypertarget{ux4ebaux7a2eux3068ux30b3ux30fcux30ebux30d0ux30c3ux30afux306eux95a2ux4fc2}{}
\begin{itemize}
\tightlist
\item
  \texttt{table()}
  関数を使って、人種（\texttt{race}）とコールバック（\texttt{call}）のクロス集計表を作成します。
\end{itemize}

\scriptsize

\begin{Shaded}
\begin{Highlighting}[]
\CommentTok{\# 人種とコールバックの2変数のクロス集計表を作成}
\CommentTok{\# race="black" または "white"}
\CommentTok{\# call=1 (連絡あり), call=0 (連絡なし)}
\NormalTok{race.call.tab }\OtherTok{\textless{}{-}} \FunctionTok{table}\NormalTok{(}\AttributeTok{race =}\NormalTok{ resume}\SpecialCharTok{$}\NormalTok{race, }\AttributeTok{call =}\NormalTok{ resume}\SpecialCharTok{$}\NormalTok{call) }

\CommentTok{\# 作成した表を表示}
\NormalTok{race.call.tab}
\end{Highlighting}
\end{Shaded}

\begin{verbatim}
##        call
## race       0    1
##   black 2278  157
##   white 2200  235
\end{verbatim}

\normalsize
\end{frame}

\hypertarget{ux30b3ux30fcux30ebux30d0ux30c3ux30afux7387ux306eux8a08ux7b97}{%
\section{2.1.3
コールバック率の計算}\label{ux30b3ux30fcux30ebux30d0ux30c3ux30afux7387ux306eux8a08ux7b97}}

\begin{frame}[fragile]{人種別のコールバック率}
\protect\hypertarget{ux4ebaux7a2eux5225ux306eux30b3ux30fcux30ebux30d0ux30c3ux30afux7387}{}
\begin{itemize}
\tightlist
\item
  集計表から、特定の人種のコールバック率を計算します。
\end{itemize}

\scriptsize

\begin{Shaded}
\begin{Highlighting}[]
\CommentTok{\# 全体のコールバック率: 連絡があった合計 / 全体の履歴書数}
\FunctionTok{sum}\NormalTok{(race.call.tab[, }\DecValTok{2}\NormalTok{]) }\SpecialCharTok{/} \FunctionTok{nrow}\NormalTok{(resume)}
\end{Highlighting}
\end{Shaded}

\begin{verbatim}
## [1] 0.08049281
\end{verbatim}

\begin{Shaded}
\begin{Highlighting}[]
\CommentTok{\# 黒人のコールバック率: 黒人で連絡あり / 黒人の合計}
\CommentTok{\# [1, 2] は「1行目(black), 2列目(call=1)」を意味する}
\NormalTok{race.call.tab[}\DecValTok{1}\NormalTok{, }\DecValTok{2}\NormalTok{] }\SpecialCharTok{/} \FunctionTok{sum}\NormalTok{(race.call.tab[}\DecValTok{1}\NormalTok{, ]) }
\end{Highlighting}
\end{Shaded}

\begin{verbatim}
## [1] 0.06447639
\end{verbatim}

\begin{Shaded}
\begin{Highlighting}[]
\CommentTok{\# 白人のコールバック率: 白人で連絡あり / 白人の合計}
\CommentTok{\# [2, 2] は「2行目(white), 2列目(call=1)」を意味する}
\NormalTok{race.call.tab[}\DecValTok{2}\NormalTok{, }\DecValTok{2}\NormalTok{] }\SpecialCharTok{/} \FunctionTok{sum}\NormalTok{(race.call.tab[}\DecValTok{2}\NormalTok{, ]) }
\end{Highlighting}
\end{Shaded}

\begin{verbatim}
## [1] 0.09650924
\end{verbatim}

\normalsize
\end{frame}

\hypertarget{ux307eux3068ux3081}{%
\section{2.1.4 まとめ}\label{ux307eux3068ux3081}}

\begin{frame}[fragile]{このセクションのまとめ}
\protect\hypertarget{ux3053ux306eux30bbux30afux30b7ux30e7ux30f3ux306eux307eux3068ux3081}{}
\begin{itemize}
\tightlist
\item
  \textbf{無作為化比較試験 (RCT)}:
  処置を無作為に割り当てることで、因果関係を科学的に検証する。
\item
  \textbf{履歴書実験}:
  名前の響き（人種）を無作為に変えることで、労働市場における差別の存在を証明した。
\item
  \textbf{Rの操作}:

  \begin{itemize}
  \tightlist
  \item
    \texttt{read.csv()} によるデータの読み込み。
  \item
    \texttt{table()} を使ったクロス集計。
  \item
    行列や表のインデックス \texttt{{[}行,\ 列{]}} による要素の抽出。
  \item
    計算結果に基づく割合（確率）の算出。
  \end{itemize}
\end{itemize}
\end{frame}

\end{document}
